\documentclass[11pt,a4paper]{article}
\usepackage[T1]{fontenc}
\usepackage[utf8]{inputenc}
\usepackage{mathptmx}
\usepackage{helvet,courier}
\usepackage[margin=24mm,headheight=15pt]{geometry}
\usepackage{amsmath,amssymb,graphicx,booktabs,array,longtable}
\usepackage[table]{xcolor}
\usepackage{microtype}
\usepackage{caption}
\usepackage{fancyhdr}
\usepackage[numbers,sort&compress]{natbib}
\ifdefined\submissionformat
\usepackage{setspace}
\usepackage{lineno}
\doublespacing
\AtBeginDocument{\linenumbers}
\fi
\usepackage{hyperref}
\input{glyphtounicode}
\pdfgentounicode=1
\definecolor{navy}{HTML}{183B56}
\definecolor{teal}{HTML}{167D8D}
\hypersetup{colorlinks=true,linkcolor=navy,citecolor=navy,urlcolor=teal,pdfauthor={},pdftitle={Graph-based exploration of chemical-disease associations}}
\captionsetup{font=small,labelfont=bf,labelsep=period}
\setlength{\parskip}{4pt}
\setlength{\parindent}{1em}
\setlength{\emergencystretch}{2em}
\renewcommand{\topfraction}{0.95}
\renewcommand{\bottomfraction}{0.9}
\renewcommand{\textfraction}{0.05}
\renewcommand{\floatpagefraction}{0.85}
\setcounter{topnumber}{3}
\setcounter{totalnumber}{5}
\renewcommand{\arraystretch}{1.18}
\pagestyle{fancy}
\fancyhf{}
\fancyhead[L]{\small\sffamily Chemical-disease association analysis}
\fancyhead[R]{\small\sffamily Draft for author review}
\fancyfoot[C]{\small\thepage}
\renewcommand{\headrulewidth}{0.3pt}
\newcommand{\blankfield}[1]{\noindent\textbf{#1}\ \rule{0.68\linewidth}{0.3pt}\par}
\newcommand{\sourceitem}[1]{\textsf{#1}}
\graphicspath{{../figures/}}

\renewcommand{\thetable}{S\arabic{table}}
\renewcommand{\thefigure}{S\arabic{figure}}
\renewcommand{\thesection}{S\arabic{section}}
\fancyhead[L]{\small\sffamily Supplementary information}
\begin{document}
\begin{center}
{\LARGE\bfseries Supplementary information}\par\vspace{8pt}
{\large Graph-based exploration of chemical--disease associations with literature and gene-context evidence}\par
\end{center}
\noindent This document records the evidence basis of the exploratory manuscript and specifies a prospective evaluation design. Reported historical outcomes, direct descriptive reanalysis and future experiments are identified separately. No model training, new literature-screening experiment or new gene-path computation was performed in preparing this analysis.

\section{Source provenance}
Seven presentation files, the interface specification, model sources, two training logs, four weight files, prediction exports and the shared-gene plotting script were inspected. Presentation text was extracted slide by slide; embedded literature charts and the gene-overlap chart were also inspected visually. Table~\ref{tab:provenance} maps the main quantitative claims to these materials. The original filenames and content hashes are retained in the accompanying \texttt{evidence/source\_inventory.json}; aliases avoid ambiguity among similarly titled progress reports.

\begin{table}[htbp]
\centering\small
\caption{Provenance of the quantitative claims. Reported results have not been reproduced from original training or screening inputs.}
\label{tab:provenance}
\begin{tabular}{p{0.28\linewidth}p{0.27\linewidth}p{0.34\linewidth}}
\toprule
Claim & Source locator & Status and interpretation \\
\midrule
CTD 9,546,982 records; direct subset 108,217 & P1, slides 2 and 4 & Reported inventory; not unique-edge counts \\
Exposure 3,125 records & P1, slides 2 and 10; P5, slide 2 & Reported literature-extracted collection \\
Exposure graph: 639 chemicals, 94 diseases & P2, slide 2 & Reported pilot graph \\
GCN/MLP/GAT exposure metrics & P2, slide 2 & Single reported table; split and averaging details unavailable \\
Seven-window literature yields & P1, slide 14 chart & Transcribed chart labels; individual judgments unavailable \\
304/1,000 and 72/1,000 literature yields & P1, slide 15 chart & Reported ranked/random retrieval outcomes \\
62,042 records; 61,973 name pairs & E1, entire CSV & Recomputed directly from retained records \\
128 shared-gene records; 639 mappable & P3, slide 1 & Reported gene screen \\
12.8\% observed; 18.7\% permutation mean & E3 chart; E4 script & Chart-reported comparison, not a rerun \\
Three case publications & P4, slide 3 hyperlinks & Bibliographic records and abstracts independently checked through PubMed \\
\bottomrule
\end{tabular}
\end{table}

P1 denotes the 15-slide disease--pollutant overview; P2 the environmental group meeting presentation dated 15 December 2025; P3 the 0809 weekly progress presentation; P4 the environmental group meeting dated 25 January 2026; and P5 the corresponding 2 February 2026 presentation. P5, slide 3 reverses the labels for the CTD and exposure counts. The main manuscript uses the assignment corroborated by P1 and P5, slide 2.

E1 is \nolinkurl{graphsage_wrong_predictions_all.csv}; E2 is \nolinkurl{graphsage_top10_wrong_predictions.csv}; E3 is \nolinkurl{mechanism_support_0809.png}; and E4 is \nolinkurl{build_weekly_progress_0809.py}. The inference-score threshold table in P1 is not used as evidence for unique graph sizes: its nested chemical counts are inconsistent, and record totals exceed possible unique-pair totals for one reported subset. This exclusion does not affect the directly counted prediction export.

\section{Model configuration and reproducibility limits}
Table~\ref{tab:config} describes the retained GraphSAGE source. Training and validation are full-graph encodings of the same training-positive edge set, followed by pair classification on each split. No explicit neighborhood sampler is used in this GraphSAGE entry point. Although GraphSAGE supports inductive applications, this implementation uses learned node identities and does not supply features for an unseen chemical or disease.

\begin{table}[htbp]
\centering\small
\caption{Retained GraphSAGE configuration. This is a source inspection, not a reconstructed run manifest.}
\label{tab:config}
\begin{tabular}{p{0.43\linewidth}p{0.46\linewidth}}
\toprule
Item & Value or behavior \\
\midrule
Node embedding & 128 dimensions; learned per node \\
Encoder & Two SAGEConv layers: 128$\to$256$\to$128 \\
Decoder & Concatenation; Linear 256$\to$128; ReLU; dropout; Linear 128$\to$2 \\
Dropout & 0.1 \\
Optimizer & Adam; learning rate 0.005; weight decay $10^{-4}$ \\
Epochs / checkpoint criterion & 80 / highest validation macro-F1 \\
Positive / negative record ratio & 1:1; unobserved negatives sampled with replacement \\
Split & Approximately 80/10/10, stratified by label at record level \\
Loader random seed & NumPy seed 42; not a complete training-seed specification \\
Graph direction & Chemical to disease; no explicit reverse edges \\
Export rule & Misclassified test records; score is maximum predicted-class softmax \\
\bottomrule
\end{tabular}
\end{table}

The loader returns five values. Six retained baseline entry points unpack four; the GCN source also fails static parsing because of an indentation error. The two available GAT logs terminate with a missing \texttt{torch\_geometric} dependency and contain no numerical epoch metrics. The original CTD and environmental input paths are absent. These observations prevent the current source collection from serving as a single verified reproduction environment. They do not establish that the earlier experiments were never run.

Weight-file metadata further indicates different task configurations: the retained GraphSAGE embedding has 25,018 nodes and a two-class output, the final GAT embedding has 728 nodes and two classes, and two other GAT weights have 13,942 nodes and three classes. These metadata were inspected without executing the serialized models. Weight files with different node counts and output classes cannot be treated as matched model-comparison results without their own input mappings and run manifests.



\section{Literature-screening aggregates}
Table~\ref{tab:windows} contains the expanded-window rates. Counts in the third column are reconstructed from the printed percentages and a stated denominator of 1,000; they are not a substitute for individual literature judgments. For the first six windows, support decreases as numerical rank increases. Thus a Spearman coefficient computed against numerical rank would be negative; a positive sign would refer to a decreasing score axis. The manuscript omits the ambiguous correlation annotation in the source chart and reports the observed sequence directly.

\begin{table}[htbp]
\centering\small
\caption{Literature-retrieval summaries from the seven-window chart.}
\label{tab:windows}
\begin{tabular}{rrrr}
\toprule
Ranks & Records examined & Reconstructed supported count & Reported yield \\
\midrule
1--1,000 & 1,000 & 304 & 30.4\% \\
10,001--11,000 & 1,000 & 193 & 19.3\% \\
20,001--21,000 & 1,000 & 158 & 15.8\% \\
30,001--31,000 & 1,000 & 155 & 15.5\% \\
40,001--41,000 & 1,000 & 145 & 14.5\% \\
50,001--51,000 & 1,000 & 144 & 14.4\% \\
60,001--61,000 & 1,000 & 178 & 17.8\% \\
\bottomrule
\end{tabular}
\end{table}

Within the first window, the ten consecutive 100-record bins have reported yields of 38, 34, 37, 30, 28, 33, 21, 31, 29 and 23\%. The corresponding printed-bin counts sum to 304, consistent with the aggregate 30.4\% yield. The original 100-per-window screening and the later 1,000-per-window screening are not pooled as independent samples because their overlap is not documented.

The uniform random-pair comparator reports 72 supported pairs among 1,000. Its vocabulary sizes, 16,171 chemical names and 6,193 disease names, match the distinct names directly counted in E1. That numerical correspondence does not establish the identity of the historical sampled pairs, deduplication policy or retrieval run. The ratio $30.4/7.2=4.222\ldots$ is therefore a descriptive comparison of the reported yields.

\section{Prediction-export analysis and gene denominators}
The reproducible descriptive script parses CSV rather than counting physical lines, treats names as strings, and checks distinct pairs and complete-row duplicates. It returns 62,042 records in E1. The source file is identified by SHA256:
\begin{center}\small\ttfamily
56f8adbe89388cb2b277665e6189be433\par
d62f8f5e38f7cd65bbbd1fb512f0efd
\end{center}
The file is sorted by decreasing score. All 69 repeated pair occurrences are exact duplicates in the stored fields. The absence of labels, stable identifiers and run identity means that an exact name-pair duplicate cannot be traced to a duplicate source record, a repeated negative draw or a later export operation. The separate E2 file is not a subset of E1 by chemical--disease name pair and is not combined with it.

\begin{table}[htbp]
\centering\small
\caption{Additional descriptive quantities and gene-screen denominators. CSV quantities are recomputed; gene-screen quantities are reported aggregates.}
\label{tab:additional}
\begin{tabular}{p{0.57\linewidth}r}
\toprule
Quantity & Value \\
\midrule
E1 mean maximum-class score & 0.7483805 \\
E1 median maximum-class score & 0.7536072 \\
E1 records with score $\ge 0.90$ & 12,413 \\
E1 records with score $\ge 0.95$ & 4,844 \\
E1 records with score $\ge 0.99$ & 200 \\
Gene-screen selected records & 1,000 \\
Mappable on both chemical and disease sides & 639 (63.9\%) \\
Shared-gene records & 128 \\
Shared-gene rate among mappable records & 20.03\% \\
Shared-gene rate among all selected records & 12.8\% \\
Mean overall rate under 200 disease-label permutations & 18.7\% \\
\bottomrule
\end{tabular}
\end{table}

For a permutation $\pi_b$ of the 1,000 disease labels, the comparator is
\[
 R_b=\frac{1}{1000}\sum_{i=1}^{1000}P(c_i,d_{\pi_b(i)}),\qquad \bar R=\frac{1}{200}\sum_{b=1}^{200}R_b.
\]
The rate among mappable original pairs is a different estimand from $R_b$. Permuting labels can alter which chemical--disease pairings have mapped values on both sides, even when each marginal multiset is preserved. Without replicate-level rates and mapping indicators, neither a conditional null comparison nor a significance calculation is available. The figure uses the matching all-record denominator only.

\section{Evidence semantics and example publications}
PMIDs 12861854, 9538487 and 1678802 were checked through the official NCBI PubMed efetch endpoint on 6 October 2026. Verification covered bibliographic metadata and abstracts, not an independent reanalysis of patient data or a full systematic review. The first describes sulfamethoxazole-associated acute generalized exanthematous pustulosis, the second a hypothesized indirect lithium-related clinical pathway, and the third a therapeutic octreotide study. These examples demonstrate why co-retrieval of two entity names cannot determine relation direction.

The evidence record proposed for a future system includes: chemical and disease identifiers; original names; relation type; model score and score class; graph snapshot and run identifier; source PMID and DOI; resolved Acemap identifier; original evidence passage and its location; organism and study design; exposure context; evidence category; reviewer status; and retrieval date. A source link, a resolved bibliographic match and scientific support for the asserted relation should be separate fields.

\section{Prospective evaluation specification}
This section is a proposed research protocol. It supplies no completed experiment results.

\textbf{Dataset freezing and identity.} Save source release/date, license, checksum and extraction criteria. Resolve chemicals with explicit namespaces and retain mixtures or classes without forcing a molecular identity. Resolve diseases to supported ontology identifiers. Preserve source-level rows but construct model splits from unique identifier pairs and typed relations. Record uncertain mappings and report their coverage.

\textbf{Labels and candidate universe.} Separate therapeutic, adverse-event, exposure-association, marker/mechanism and inferred relationships. Use a declared binary task only when the collapsed meaning is appropriate. Label absent pairs as unobserved; document negative sampling and evaluate sensitivity to its choice. Export all positive-class candidate scores outside the frozen known-pair set, with $y$, $\hat y$, score vector, split, entity IDs, pair ID, snapshot and run ID.

\textbf{Splits and leakage checks.} Publish unique-pair train/validation/test manifests. Require zero held-out pair overlap with training labels and message-passing edges, including reverse-edge and source-alias checks. Where gene-mediated edges derive from target associations, check target-dependent inference leakage. Assess transductive random, chemical-disjoint, disease-disjoint and temporal splits as distinct settings. Specify how unseen nodes are represented before interpreting disjoint evaluations.

\textbf{Comparators and metrics.} Compare degree-based scores, an embedding MLP and graph baselines on identical splits. Include a disease-association-specific method such as LAGCN or MRDDA when its inputs and implementation can be made comparable. Predefine tuning budgets and at least five independently seeded runs as a practical project target. Report AUROC, average precision, precision at fixed review budgets and calibration. Use paired comparisons and uncertainty estimates appropriate to repeated entities; disclose how prevalence in the test set differs from the application candidate space.

\textbf{Literature validation.} Freeze the prediction list before searching. Record query strings, date, retrieval system, model version if used, and retained bibliographic identifiers. Use an explicit rubric: direct relation support; indirect/contextual support; therapeutic relation; contradictory result; co-mention only; or no eligible evidence. Obtain independent judgments and adjudicate disagreement. Separate pre-cutoff literature from later evidence and remove articles used to construct training labels from any independent-validation count. Compare uniform and degree-/literature-matched controls using the same retrieval budget.

\textbf{Gene contribution.} Compare chemical--disease-only models with typed gene-augmented models, a shuffled-gene control and degree-aware baselines. Report mapping coverage and shared-path rate separately. Keep each path's two source edges and evidence IDs. Evaluate whether gene information improves held-out prediction and whether it yields useful explanations as different endpoints. Repeat the disease-label permutation and save every replicate rate before drawing enrichment conclusions.

\textbf{System evaluation.} Select a reproducible set of queries and candidate records. Measure exact entity resolution, article-link resolution and source-passage correctness. Check PubMed/DOI and Acemap links independently. Evaluate representative end-to-end user tasks and report the tested version. Availability of a homepage is insufficient evidence of a functioning literature or prediction service.

\section{System scope and remaining extensions}
The system requirements can be organized into four independently verifiable outcomes (Table~\ref{tab:scope}). Current artifacts support a broad CTD inventory, an association baseline, evidence-screening summaries and interface designs. They do not establish an integrated, evaluated release of all four outcomes.

\begin{table}[htbp]
\centering\small
\caption{System extension specification. These are intended outcomes, not completed features claimed by the manuscript.}
\label{tab:scope}
\begin{tabular}{p{0.23\linewidth}p{0.30\linewidth}p{0.35\linewidth}}
\toprule
Requested extension & Current basis & Verifiable outcome \\
\midrule
Broader chemicals & CTD inventory; SIDER-format files present & Versioned imports, identity crosswalks, typed source records and coverage report \\
Closer Acemap integration & Literature fields and interface design & PMID/DOI-to-Acemap mapping; verified article links; unresolved status retained \\
Gene nodes and relations & Shared-gene screening script and chart & Typed, sourced gene paths; ablation against matched models and nulls \\
Literature for new candidates & Retrieval-rate summaries and ten case links & Positive-class candidate export; auditable evidence judgments; independent validation split \\
\bottomrule
\end{tabular}
\end{table}


\end{document}
